% Source: https://artofproblemsolving.com/community/c3h3186635_2021_abmc_speed_contest__actonboxborough_math_competition
% Stable identity: aops-topic-3186635-q23
Write $a\circ b=2a+\dfrac{(a+4)(a-3)}{b}$. For every inner level $a\ge4$, the denominator is positive and the result is positive, so all expressions are defined. At $a=3$, the second term vanishes: $3\circ b=6$. Consequently the entire expression is $1\circ(2\circ6)=1\circ3=2-\dfrac{10}{3}=-\dfrac{4}{3}$. Directly, $2\circ6=3$, and $1\circ3=\dfrac{1+6+1-12}{3}=-\dfrac{4}{3}$. Hence $a+b=\boxed7$.
